\section{Valeurs propre}
\subsection{Valeurs propres et vecteurs propres}

\newdef{Soit $V$ un $\K$-espace vectoriel et $f : V \to V$ un endomorphisme. Un \emph{vecteur propre} associé à la \emph{valeur propre} $\lambda \in \K$ est un vecteur $v \in V^*$ tel que
	\[
	f(v) = \lambda v
	\]
}{}

\newlem{Soit $B = \{v_1, \ldots, v_n\}$ une base de $V$ et $A \in M_n(\K)$ la matrice de l'endomorphisme $f : V \to V$ relatif à $B$. La matrice $A$ est diagonale, c'est-à-dire $A$ de la forme
	\[
	A = \begin{pmatrix}
		\lambda_1 & & \\
		          & \ddots & \\
		          &        & \lambda_n
	\end{pmatrix}
	\]
	si et seulement si $v_i$ est un vecteur propre associé à la valeur propre $\lambda_i$ pour tout $i \in \entiers$.
}{}

\prf{Supposons que $v_i$ soit un vecteur propre associé à la valeur propre $\lambda_i \in \K$, ainsi dans la base $B$ on a
	\[
	[f(v_i)]_B = \lambda [v_i]_B = \lambda e_i
	\]
	ainsi on a bien,
	\[
	A = \begin{pmatrix}
		\lambda_1 & & \\
		& \ddots & \\
		&        & \lambda_n
	\end{pmatrix}
	\]
La réciproque est analogue.
}

\newdef{Soit $\lambda \in \K$ une valeur propre de l'endomorphisme $f : V \to V$, alors
	\[
	E_{\lambda} = \ker(f - \lambda \mathrm{Id})
	\]
	est appelé l'\emph{espace propre} de $\lambda$. De plus $\dim E_\lambda$ est appelé \emph{multiplicité géométrique} de $\lambda$.
}{}

\newlem{Soit $f : V \to V$ un endomorphisme et $v_1, \ldots, v_k \in V$ des vecteurs propres associés aux valeurs propre $\lambda_1, \ldots, \lambda_k \in \K$ deux à deux distincts. Alors le famille $\{v_1, \ldots, v_k\}$ est libre.
}{}

\prf{Par récurrence sur $k$, pour $k = 1$ comme $v_1 \neq 0$, alors $\{v_1\}$ est libre. Pour $k > 1$, on considère une combinaison linéaire, sans perte de généralité supposons que $\lambda_k \neq 0$, 
	\[
	\alpha_1 v_1 + \ldots \alpha_k v_k = 0
	\]
Par évaluation de cette égalité par l'endomorphisme $f$ et division par $\lambda_k$ on a
\[
\alpha_1 \frac{\lambda_1}{\lambda_k} v_1 + \ldots + \alpha_k v_k = 0
\]
ainsi en soustrayant les deux équations on obtient
\[
\alpha_1 \left(1 - \frac{\lambda_1}{\lambda_k}\right) v_1 + \ldots + \alpha_{k-1} \left( 1 - \frac{\lambda_{k-1}}{\lambda_k}\right) v_{k-1} = 0
\]
ainsi par récurrence on en déduit que,
\[
\forall i \in \llbracket 1, k-1 \rrbracket, \quad \alpha_i \underbrace{\left(1 - \frac{\lambda_i}{\lambda_k}\right)}_{\neq 0} = 0 \implies \forall i \in \llbracket 1, k-1 \rrbracket, \quad \alpha_i = 0
\]
}

\newlem{Soit $f : V \to V$ un endomorphisme et $\dim V = n$. Soient $\lambda_1, \ldots, \lambda_r$ les valeurs propres deux à deux distinctes. On considère
	\[
	B_i = \{ v_1^i, \ldots, v_{n_i}^i\}
	\]
	les bases respectives de $E_{\lambda_i}$. Alors
	\[
	B = \bigcup_{i=1}^r B_i
	\]
	est un ensemble libre. De plus,
	\[
	f \text{ diagonalisable} \iff n_1 + n_2 + \ldots + n_r = n
	\]
}{Base des espaces propres}

\prf{
En premier temps, on considère une combinaison linéaire tel que,
\[
\sum_{i = 1}^r \underbrace{\sum_{j = 1}^{n_i} \alpha_{ij} v_j^i}_{\in E_{\lambda_i}} = 0
\]
mais tenant compte du lemme précédent, on en déduit que
\[
\sum_{j = 1}^{n_i} \alpha_{ij} v_j^i = 0
\]
mais comme $B_i$ est une base, en particulier c'est une famille libre, ainsi,
\[
\forall i,j \quad \alpha_{ij} = 0
\]
Donc $B$ est bien un famille libre. De plus si on suppose que $\sum n_i = n$, alors $B$ est un base, mais comme $B$ est une base de vecteur propre, alors par définition  $f$ est diagonalisable. Réciproquement si $f$ est supposé diagonalisable, et si $m_i$ dénote le nombre de vecteur propre en $E_{\lambda_i}$ dans la base relatif consistant de vecteur propre, alors $m_i \leqslant n_i$, et donc
\[
n = m_1 + \ldots + m_r \leqslant n_1 + \ldots + n_r \leqslant n
\]
et donc $n_1 + \ldots + n_r = n$.
}


\subsection{Polynôme caractéristique}

\newdef{Soit $A$ la matrice associé à l'endomorphisme de $f$, le \emph{polynôme caractéristique} de $f$ est donnée par,
	\[
	\chi_A(x) = \det (A - x I_n) = \sum_{\sigma \in S_n} \mathrm{sign}(\sigma) \prod_{i = 1}^n (A- x I_n)_{i \sigma(i)} = \alpha_0 + \alpha_1 x + \ldots + \alpha_n x^n
	\]
	où $\alpha_i \in \K$. Pour $\lambda$ une valeur propre de l'endomorphisme $f$, la \emph{multiplicité algébrique} de $\lambda$ est la multiplicité de $\lambda$ comme racine de $\det(A - x I_n)$.
}{Polynôme caractéristique}

\newprop{Soit $f : V \to V$ un endomorphisme et soit $\lambda \in \K$ une valeur propre $f$. Alors la multiplicité géométrique est au plus la multiplicité algébrique.
}{Majoration de la multiplicité géométrique}

\prf{Soit $\lambda \in \K$ une valeur propre de $f$ et $B_\lambda = \{ v_1, \ldots, v_k\}$ une base de 
	\[
	E_\lambda = \ker(f - \lambda \mathrm{Id})
	\]
	On complète alors $B_\lambda$ est une base $B'$ de $V$,
	\[
	B' = \{v_1, \ldots, v_k, u_1, \ldots, u_{n-k}\}
	\]
	Ainsi la matrice associé à notre endomorphisme dans la base $B'$ est,
	\[
	A = \begin{pmatrix}
		\lambda I_k & D \\
				0 & C
	\end{pmatrix}
	\]
	ainsi le polynôme caractéristique associé est,
	\[
	\chi_A (x) = (\lambda - x)^k \det(C - x \lambda I_{n-k})
	\]
	et donc 
	\[
	(\lambda- x)^k ~|~\chi_A(x)
	\]
	donc la multiplicité de $\lambda$ est d'au moins $k$.
} 

\newthm{Soit $V$ un espace vectoriel de dimension $n$, $f : V \to V$ un endomorphisme et $\lambda_1, \ldots, \lambda_k \in \K$ les valeurs propres distincts de $f$. Alors $f$ est diagonalisable si et seulement si
	\begin{enumerate}
		\item Le polynôme caractéristique $\chi_f$ de $f$ est scindé,
		\[
		\chi_f(x) = (-1)^n \prod_{i =1}^k (x - \lambda_i)^{\alpha_i}
		\]
		où $\alpha_i$ est la multiplicité algébrique de $\lambda \in K$.
		\item Les multiplicités algébriques et géométrique sont les mêmes.
	\end{enumerate}
}{}

\prf{Supposons $f$ diagonalisable. D'après le lemme (\ref{Base des espaces propres}) alors,
	\[
	\dim E_{\lambda_1} + \ldots + \dim E _{\lambda_k} = n
	\]
	mais d'après la proposition (\ref{Majoration de la multiplicité géométrique})
	\[
	\alpha_1 + \ldots + \alpha_k \leqslant n
	\]
	ainsi on en déduit que,
	\[
	\forall i, \quad \dim E_{\lambda_i} = \alpha_i
	\]
	Donc la multiplicité géométrique et algébrique coïncident bien pour toutes les valeurs propres. Réciproquement, d'après le polynôme caractéristique on en déduit que,
	\[
	\alpha_1 + \ldots + \alpha_k = n
	\]
	mais d'après la deuxième hypothèse on a
	\[
	\forall i, \quad \dim E_{\lambda_i} = \alpha_i
	\]
	on en conclut par le lemme (\ref{Base des espaces propres}) que $f$ est diagonalisable.
}

\subsection{Théorème de Cayley-Hamilton}

\newthm{Soit $A \in M_n(\K)$ et $\chi_A$ le polynôme caractéristique de $A$, alors
	\[
	\chi_A(A) = 0
	\]
}{Théorème de Cayley-Hamilton}

\prf{Pour $A$ un matrice diagonalisable alors 
	\[
	A = P D P^{-1}
	\]
	de plus si $A$ est diagonalisable alors $\chi_A = (-1)^n \prod_{i = 1}^n (x - \lambda_i)$ et donc 
	\[
	\chi_A(A) = \chi_A(P D P^{-1}) = P \chi_A(D) P^{-1} = P \cdot 0 \cdot P^{-1}
	\]
	ainsi $\chi_A(A) = 0$ si $A$ est diagonalisable. Et donc par densité des matrices diagonalisables \footnote{Cette démonstration utilisant la densité des matices diagonalisables n'est \textbf{ABSOLUMENT} pas au programme (de Algèbre linéaire 2) mais je trouvais cette démonstration très jolie. Donc si tu lis ça, tu seras peut-être content de savoir que les matrices diagonalisables sont denses dans $M_n(\K)$ au même titre que les matrices inversibles.} dans $M_n(\K)$ et par continuité de $M \mapsto \chi_M(M)$, alors $\chi_A(A) = 0$ pour une matrice quelconque.
}